\vspace{-2pt}
\section{Method}
\vspace{-3pt}
% This section describes EPA, which involves instructing an LLM to add general ASR errors to the overall sentence, followed by positioning errors to keywords using in-context learning (Figure ~\ref{fig:method}).

\subsection{Notation}
Before detailing each step, we first clarify the notation. Dialogue context from turn 1 to $t$ is denoted as $D_t$=$\{(s_1, u_1 ),..., ( s_t, u_t )\}$ where $s$ denotes for system and $u$ for user utterance. DST model predicts the dialogue state (also called belief state) $B_t$ given $D_t$. $B_t$ is composed with slot $sl$ and value $v$  pairs, denoted as $B_t$ = $\{( {sl}_1, v_1 ),..., ( {sl}_J, v_J )\}$ , where ${sl}_j$ and $v_j$ is $j$-th slot name and value. $J$ is the total number of slots.

\subsection{Step 1: ASR Error for Overall Utterance}
\label{sec:method_step1}
In this step, we augmented the overall utterance by introducing general ASR errors. We began by constructing example sets for in-context learning, utilizing an open-source audio dataset \cite{ardila-etal-2020-common}. From this dataset, we randomly selected 300 hours of audio along with their corresponding gold transcripts ($g$) and transcribed the audio using an off-the-shelf ASR model (e.g., Whisper-base \cite{whisper}) to obtain the erroneous transcriptions ($e$). We denote this example dataset as $DB = \{(g_1, e_1), \ldots, (g_I, e_I)\}$.


Next, we inject errors into $u$ by prompting the LLM with in-context examples. We retrieved ($g$, $e$) pairs from the database ($DB$) based on phonetic similarity between $u$ and $g$ (Figure~\ref{fig:method}, Step 1). To compute phonetic similarity, we converted the characters of both $u$ and $g$ into phonemes using the International Phonetic Alphabet (IPA), and calculated similarity using a frequency-based retrieval \footnote{We used BM25\cite{bm25}, a retrieval model based on term frequency. While neural retrievers (e.g., DPR\cite{DPR}) could be applied, we opted for a frequency-based method, as neural models tend to capture semantic similarity.}. After selecting the top-$k$ ($g$, $e$) pairs, we concatenated the instruction, in-context examples, and $u$ into a single prompt and provided it to the LLM. This process results in the overall ASR-errored user utterance, denoted as $\dot{u}$. Concretely, $\dot{u}$ can be obtained by 
 \begin{equation}
    \dot{u_t} = LLM(Inst_1 \oplus (g_1, e_1)\cdots(g_k, e_k)\oplus u_t)
\end{equation}
where $\oplus$ denotes concatenation, and we set $k=3$ throughout our experiments.  Retrieved examples are provided in Appendix~\ref{app:EPA_retreive}.

\subsection{Step 2: ASR Error for Keywords}
\label{sec:method_step2}



% While Step 1 introduces general ASR errors into $u$, it does not guarantee increased error diversity in keywords. To construct a more effective training dataset, we explicitly generated keyword-focused ASR errors (Figure~\ref{fig:method}, Step 2). For this process, we provided in-context examples that highlight the keywords in the clean text using the <hl> tag ($\tilde{g}$), and pair them with the corresponding ASR-errored versions ($\tilde{e}$)\footnote{For the DST task, we treated \textit{values ($v$)} for dialogue state—such as hotel or restaurant names—as keywords, although the definition of a keyword may differ by task.}.We applied the same highlighting to the inference target $\dot{u}$ and then concatenated the instructions, examples, and highlighted $\dot{u}$ into a single prompt.
% Using this prompt, the model generates the augmented utterance $\ddot{u}$ with both overall and keyword-specific ASR errors. Explicitly, we get $\ddot{u}$ by
% \begin{equation}
%     \ddot{u_t} = LLM(Inst_2 \oplus (\tilde{g}_0, \tilde{e}_0)\cdots(\tilde{g}_k, \tilde{e}_k)\oplus \dot{u}_t).
% \end{equation}
% The used prompts are provided in Appendix~\ref{app:prompt}.

% \subsection{Examples of EPA}
% 
\begin{table}[t]
\resizebox{\columnwidth}{!}{%
\begin{tabular}{c|l|l}
\hline
            Idx &    Method & \multicolumn{1}{c}{Examples}                           \\ \hline
\multirow{2}{*}{1} & Original & Tuesday, going to bailey's crossroads please.           \\
                   & +EPA & Tuesday, going to \colorbox{orange}{baley's} crossroads, \colorbox{yellow(ryb)}{peas}.             \\ \hline            
\multirow{2}{*}{1} & Original & I'd like to find a vegetarian restaurant, if possible.  \\
                   & +EPA & I'd \colorbox{yellow(ryb)}{hike} to find a \colorbox{orange}{veggie tarian} \colorbox{yellow(ryb)}{restroom}, if possible. \\ \hline
\multirow{2}{*}{3} & Original & I am going to auburn.                                   \\
                   & +EPA & I am \colorbox{yellow(ryb)}{flowing} to \colorbox{orange}{auburng}.                                \\ \hline

\multirow{2}{*}{4} & Original & Hi! Could you find me a train to loris on thursday?     \\
                   & +EPA & \colorbox{yellow(ryb)}{Oh}! Could you find me a \colorbox{yellow(ryb)}{trai} to \colorbox{orange}{lorri} on \colorbox{orange}{thursdae}?      \\ \hline
\multirow{2}{*}{5} & Original & Ashby is my destination.                                \\
                   & +EPA & \colorbox{orange}{Ashy's} my \colorbox{yellow(ryb)}{desity}.                                       \\ \hline
\end{tabular}%
}
\caption{Examples of ASR errors from EPA. }
\label{tab:EPA_examples}
\vspace{-15pt}
\end{table}

% Table~\ref{tab:EPA_examples} shows examples of ASR errors generated by EPA. We have highlighted utterance level overall errors in \colorbox{yellow(ryb)}{yellow} and keyword-specific errors in \colorbox{aqua}{blue}. For instance, in Row 1, the model introduces a keyword-level error (bailey’s → baley’s) as well as an additional phonetically plausible insertion (peas), simulating realistic ASR noise. Further examples can be found in Appendix~\ref{app:asr-examples}.




While Step 1 introduces general ASR-style errors into $u$, it does not ensure sufficient error diversity in keyword tokens. To construct a more effective training dataset, we explicitly generate keyword-focused ASR errors in Step 2 (Figure~\ref{fig:method}). In this step, we highlight the keywords in $\dot{u}$ using the \texttt{<hl>} tag and instruct the LLM to inject errors specifically within the highlighted spans. For the DST task, we treat dialogue state values ($v$) as keywords, although the definition of a keyword may vary depending on the task. To facilitate this process, we provide a few examples that illustrate how values within \texttt{<hl>} tags are intended to be modified during augmentation. Given these instructions and examples, the LLM generates an augmented utterance $\ddot{u}$ that includes both general and keyword-specific ASR errors. Formally, we obtain $\ddot{u}_t$ as follows: 
\begin{equation}
    \ddot{u_t} = LLM(Inst_2 \oplus ({g}_0, {e}_0)\cdots({g}_k, {e}_k)\oplus \dot{u}_t).
\end{equation}
The used prompts are provided in Appendix~\ref{app:prompt}.

\subsection{Examples of EPA}

\begin{table}[t]
\resizebox{\columnwidth}{!}{%
\begin{tabular}{c|l|l}
\hline
            Idx &    Method & \multicolumn{1}{c}{Examples}                           \\ \hline
\multirow{2}{*}{1} & Original & Tuesday, going to bailey's crossroads please.           \\
                   & +EPA & Tuesday, going to \colorbox{orange}{baley's} crossroads, \colorbox{yellow(ryb)}{peas}.             \\ \hline            
\multirow{2}{*}{1} & Original & I'd like to find a vegetarian restaurant, if possible.  \\
                   & +EPA & I'd \colorbox{yellow(ryb)}{hike} to find a \colorbox{orange}{veggie tarian} \colorbox{yellow(ryb)}{restroom}, if possible. \\ \hline
\multirow{2}{*}{3} & Original & I am going to auburn.                                   \\
                   & +EPA & I am \colorbox{yellow(ryb)}{flowing} to \colorbox{orange}{auburng}.                                \\ \hline

\multirow{2}{*}{4} & Original & Hi! Could you find me a train to loris on thursday?     \\
                   & +EPA & \colorbox{yellow(ryb)}{Oh}! Could you find me a \colorbox{yellow(ryb)}{trai} to \colorbox{orange}{lorri} on \colorbox{orange}{thursdae}?      \\ \hline
\multirow{2}{*}{5} & Original & Ashby is my destination.                                \\
                   & +EPA & \colorbox{orange}{Ashy's} my \colorbox{yellow(ryb)}{desity}.                                       \\ \hline
\end{tabular}%
}
\caption{Examples of ASR errors from EPA. }
\label{tab:EPA_examples}
\vspace{-15pt}
\end{table}

Table~\ref{tab:EPA_examples} shows examples of ASR errors generated by EPA. We have highlighted utterance level overall errors in \colorbox{yellow(ryb)}{yellow} and keyword-specific errors in \colorbox{orange}{orange}. For instance, in Row 1, the model introduces a keyword-level error (bailey’s → baley’s) as well as an additional phonetically plausible insertion (peas), simulating realistic ASR noise. Further examples can be found in Appendix~\ref{app:asr-examples}.


